\section{Tutorial/Functionality}
\label{sec:tutorial_advanced}

In this chapter, we will dissect the codebase developed for real-time hand detection and finger counting. Beyond simply presenting the code, this section elucidates the underlying algorithmic thought process and details the fundamental commands provided by the \textbf{OpenCV} and \textbf{MediaPipe} libraries.

\subsection{Algorithmic Rationale}
Before delving into the syntactic implementation, it is crucial to understand the logical pipeline of the system:
\begin{enumerate}
    \item \textbf{Data Acquisition:} Capture a continuous stream of visual data from a hardware peripheral (webcam).
    \item \textbf{Normalization:} Mirror the image for intuitive user interaction and convert the color space to meet the machine learning model's prerequisites.
    \item \textbf{Inference:} Pass the normalized image through the MediaPipe neural network to extract a topological graph of the hand (landmarks).
    \item \textbf{Kinematic Evaluation:} Utilize the spatial coordinates of these landmarks to deduce the state of each digit (extended or folded) based on anatomical constraints.
    \item \textbf{Feedback:} Overlay the computed state onto the original image stream to provide real-time visual feedback.
\end{enumerate}

\subsection{Environment Setup and Core Libraries}
\label{subsec:setup}

We initiate the script by importing the requisite libraries and defining the core MediaPipe objects.

\begin{lstlisting}
import os
os.environ["QT_QPA_PLATFORM"] = "xcb" 

import cv2
import mediapipe as mp

mp_hands = mp.solutions.hands
mp_drawing = mp.solutions.drawing_utils
hands = mp_hands.Hands(min_detection_confidence=0.5, min_tracking_confidence=0.5)
\end{lstlisting}

\textbf{Basic Commands Explained:}
\begin{itemize}
    \item \texttt{import cv2}: Imports the OpenCV library, the industry standard for real-time computer vision and image processing tasks.
    \item \texttt{mp.solutions.hands}: Accesses the specific hand tracking module within the broader MediaPipe framework.
    \item \texttt{mp\_hands.Hands(...)}: Instantiates the machine learning model. The parameters \texttt{min\_detection\_confidence} and \texttt{min\_tracking\_confidence} act as statistical thresholds. Setting them to $0.5$ dictates that the model must be at least 50\% certain that a hand is present to begin or continue tracking it, balancing processing speed with accuracy.
\end{itemize}

\subsection{Video Capture and Landmark Mapping}
\label{subsec:capture}

The subsequent phase involves initializing the hardware interface and establishing the indices for the finger landmarks.

\begin{lstlisting}
cap = cv2.VideoCapture(1)
cv2.namedWindow('Hand Tracker', cv2.WINDOW_NORMAL)

# IDs for the fingertip landmarks in MediaPipe (Index, Middle, Ring, Pinky)
pontas_dos_dedos = [8, 12, 16, 20]
ponta_do_polegar = 4
\end{lstlisting}

\textbf{Basic Commands Explained:}
\begin{itemize}
    \item \texttt{cv2.VideoCapture(1)}: Creates a video capture object. The argument \texttt{1} targets the secondary camera interface, while \texttt{0} typically targets the primary built-in webcam.
    \item \texttt{cv2.namedWindow(...)}: Initializes a graphical user interface (GUI) window. The \texttt{WINDOW\_NORMAL} flag allows the user to resize the output window dynamically.
\end{itemize}

\subsection{Image Pre-processing: The "Why" and "How"}
\label{subsec:preprocessing}

Within the continuous execution loop (\texttt{while True}), raw frames are acquired and mathematically transformed.

\begin{lstlisting}
while True:
    ret, frame = cap.read()
    if not ret:
        break

    frame = cv2.flip(frame, 1)
    rgb_frame = cv2.cvtColor(frame, cv2.COLOR_BGR2RGB)
    results = hands.process(rgb_frame)
\end{lstlisting}

\textbf{Thought Process \& Basic Commands:}
\begin{itemize}
    \item \texttt{cap.read()}: Extracts the subsequent frame from the video stream. It returns a boolean (\texttt{ret}) indicating success, and the matrix of pixels (\texttt{frame}).
    \item \texttt{cv2.flip(frame, 1)}: Applies a horizontal mirror reflection across the vertical axis (indicated by the \texttt{1}). Without this operation, moving your hand to the left would result in the on-screen hand moving to the right, which causes severe cognitive dissonance for the user.
    \item \texttt{cv2.cvtColor(...)}: OpenCV natively captures color arrays in Blue-Green-Red (BGR) format. However, the MediaPipe neural network was trained on Red-Green-Blue (RGB) datasets. This function strictly translates the color channels to ensure the model processes the image accurately.
    \item \texttt{hands.process(rgb\_frame)}: Executes the core inference step, passing the RGB matrix into the neural network and returning a \texttt{results} object containing the spatial coordinates of any detected hands.
\end{itemize}

\subsection{Kinematic Logic: The Coordinate System}
\label{subsec:logic}

When hands are detected, the algorithm iterates through them and evaluates the spatial configuration of the digits.

\begin{lstlisting}
    if results.multi_hand_landmarks:
        for hand_idx, hand_landmarks in enumerate(results.multi_hand_landmarks):
            mp_drawing.draw_landmarks(frame, hand_landmarks, mp_hands.HAND_CONNECTIONS)

            hand_label = results.multi_handedness[hand_idx].classification[0].label
            dedos_levantados = 0

            # 1. Thumb Logic (X-Axis)
            if hand_label == "Right":
                if hand_landmarks.landmark[ponta_do_polegar].x < hand_landmarks.landmark[ponta_do_polegar - 1].x:
                    dedos_levantados += 1
            else: # "Left"
                if hand_landmarks.landmark[ponta_do_polegar].x > hand_landmarks.landmark[ponta_do_polegar - 1].x:
                    dedos_levantados += 1

            # 2. Logic for the remaining 4 fingers (Y-Axis)
            for ponta in pontas_dos_dedos:
                if hand_landmarks.landmark[ponta].y < hand_landmarks.landmark[ponta - 2].y:
                    dedos_levantados += 1
\end{lstlisting}

\textbf{Thought Process for Finger Counting:}
The algorithm's core logic is predicated on the digital image coordinate system, where the origin $(0,0)$ is located at the top-left corner of the image. Consequently, moving downwards increases the Y-value, and moving rightwards increases the X-value.

\begin{itemize}
    \item \textbf{The Four Fingers (Index to Pinky):} To determine if a finger is "up", we compare the Y-coordinate of the fingertip (\texttt{landmark[ponta]}) to the Y-coordinate of the knuckle two joints below it (\texttt{landmark[ponta - 2]}). Due to the inverted Y-axis, a numerically smaller Y-value indicates a higher physical position. Thus, if $Y_{tip} < Y_{knuckle}$, the finger is extended.
    \item \textbf{The Thumb Anomaly:} The thumb operates on a different biomechanical plane, rotating horizontally rather than folding vertically. Therefore, we evaluate the X-axis. A right thumb pointing leftwards (towards the center of the body) has a smaller X-value than its base joint. This logic must be strictly inverted for the left hand, necessitating the \texttt{hand\_label} conditional check.
\end{itemize}

\subsection{Visual Output and Resource Deallocation}
\label{subsec:output}

Finally, the extracted data is overlaid onto the original BGR frame for the user's viewing.

\begin{lstlisting}
            cv2.putText(frame, f'{hand_label} hand: {dedos_levantados} fingers', (10, 50 + (hand_idx * 50)),
                        cv2.FONT_HERSHEY_SIMPLEX, 1, (0, 255, 0), 2, cv2.LINE_AA)

    cv2.imshow('Hand Tracker', frame)

    if cv2.waitKey(1) & 0xFF == ord('q'):
        break

cap.release()
cv2.destroyAllWindows()
\end{lstlisting}

\textbf{Basic Commands Explained:}
\begin{itemize}
    \item \texttt{cv2.putText(...)}: Renders text directly onto the pixel matrix. The algorithmic offset \texttt{50 + (hand\_idx * 50)} ensures that if multiple hands are detected, their respective text labels stack vertically without visual occlusion.
    \item \texttt{cv2.imshow(...)}: Commands the operating system to refresh the GUI window with the newly processed frame.
    \item \texttt{cv2.waitKey(1)}: Halts execution for 1 millisecond to listen for keyboard interrupts. If the user presses the 'q' key, the \texttt{while} loop breaks.
    \item \texttt{cap.release() \& cv2.destroyAllWindows()}: These cleanup functions are crucial; they sever the connection to the hardware camera to prevent system locks and safely close all active graphical windows.
\end{itemize}

\subsection{Advanced Integration: Custom AI Model and Mouse Control (HCI)}

While in Section 4.5 successfully count fingers, they are highly dependent on 2D perspective and camera angles. To make the system truly robust, we can replace these hardcoded rules with a custom Machine Learning (ML) model and utilize the predictions for Human-Computer Interaction (HCI), specifically controlling the operating system's mouse pointer.

\subsubsection{Feature Engineering (Wrist Relativity)}
To train a robust model, the input data must be position-invariant. If raw screen coordinates are fed into the model, the algorithm will overfit to the hand's pixel location on the screen rather than learning its structural shape. 

We set the wrist (Landmark 0) as the origin point $(0,0,0)$. All other 20 landmarks are calculated relative to the wrist using the following transformations:

\begin{align}
    x_{rel} &= x_i - x_{wrist} \\
    y_{rel} &= y_i - y_{wrist} \\
    z_{rel} &= z_i - z_{wrist}
\end{align}

\begin{lstlisting}[language=Python, caption=Extracting 3D Position-Invariant Features]
landmark_data = []
wrist = hand_landmarks.landmark[0]

for landmark in hand_landmarks.landmark:
    rel_x = landmark.x - wrist.x
    rel_y = landmark.y - wrist.y
    rel_z = landmark.z - wrist.z
    landmark_data.extend([rel_x, rel_y, rel_z])
\end{lstlisting}

\subsubsection{Training the Classifier}
By capturing a dataset containing roughly 100-200 variations of each gesture (0 to 5 fingers) using these engineered relative coordinates, the data becomes tabular. This allows us to utilize the \texttt{scikit-learn} library to train a Random Forest Classifier.

\begin{lstlisting}[language=Python, caption=Training a Random Forest Model]
import pandas as pd
from sklearn.model_selection import train_test_split
from sklearn.ensemble import RandomForestClassifier
import pickle

# Load dataset and split into 80% Training and 20% Testing
df = pd.read_csv('hand_data.csv')
X = df.drop('label', axis=1) 
y = df['label']              
X_train, X_test, y_train, y_test = train_test_split(X, y, test_size=0.2, random_state=42)

# Train the model
model = RandomForestClassifier(n_estimators=100, random_state=42)
model.fit(X_train, y_train)

# Save the trained model
with open('hand_model_3d.pkl', 'wb') as f:
    pickle.dump(model, f)
\end{lstlisting}

\subsubsection{Real-Time Mouse Control (HCI)}
With a trained model (\texttt{.pkl} file) capable of recognizing gestures, we can import the \texttt{pyautogui} library to bridge the gap between Python and the operating system. 

The tip of the index finger (Landmark 8) dictates the cursor's spatial $X$ and $Y$ coordinates, mapped to the monitor's resolution. When the ML model classifies the hand structure as a closed fist ("0"), the script triggers a left-click event.

\begin{lstlisting}[language=Python, caption=HCI Mouse Control Integration]
import pyautogui
import time

# Mouse mapping logic using the Index Finger
index_finger_tip = hand_landmarks.landmark[8]
mouse_x = int(index_finger_tip.x * screen_w)
mouse_y = int(index_finger_tip.y * screen_h)
pyautogui.moveTo(mouse_x, mouse_y, _pause=False)

# Clicking logic based on ML prediction
if prediction == '0':
    if (current_time - last_click_time) > 1.0: # 1-second cooldown
        pyautogui.click()
        last_click_time = current_time
\end{lstlisting}